\section{Discussion, Boundary Contracts, and Compliance Protocol}
\label{sec:discussion}

\subsection{Four Empirical Boundary Conditions for Physics-Constrained Models}
The scoping audit, algebraic code examination, and prospective ablations converge on four boundary conditions that determine when physical priors help, when they are neutral, and when they harm.

\paragraph{Sample Restriction and Prior Misspecification Vulnerability}
On Penmanshiel ($N_{\text{train}}=1,583$), the calibrated power curve prior at $\lambda=1.0$ reduces asymmetric dispatch cost from $0.09881 \pm 0.00027$ to $0.09760 \pm 0.00044$ ($-1.22\%$). A scale-matched scrambled pseudo-prior raises it to $0.16611 \pm 0.00335$ ($+68.1\%$). The divergence between these two outcomes---authentic priors securing modest downside protection while corrupted priors inflict severe operational damage---demonstrates a fundamental risk-return asymmetry: under strong regularization, whatever information the prior encodes is amplified, including its misspecifications. Cross-site comparisons cannot isolate this effect because they confound site topography with sample volume. In the data-rich regime (Kelmarsh, $N_{\text{train}}=5,098$), unconstrained DLinear ($0.08812$), genuine physics ($0.08791$), and scrambled prior ($0.08801$) are numerically clustered, indicating hypothesis-space saturation.

\paragraph{Low-Penalty Scale and Non-Specific Numerical Regularization}
At $\lambda = 0.01$ on Kelmarsh, a genuine physical prior ($0.08810 \pm 0.00009$) and a structurally scrambled pseudo-prior ($0.08801 \pm 0.00008$) yield closely matched decision costs with overlapping seed margins. When random power curve permutations reproduce the same outcome as domain aerodynamics, the penalty term behaves empirically as a non-specific regularizer / numerical smoothing term rather than encoding domain physical laws. Published claims of domain physical benefit under $\lambda \in [10^{-4}, 10^{-2}]$ require negative-control ablations and sensitivity testing at $\lambda \ge 0.1$ to be credible.

\paragraph{Monotonic Cost Inflation under Training Proxy Observation Staleness}
When physical constraint formulations rely on meteorological proxies to query empirical power curves during loss computation, practitioners may be tempted to substitute lagged historical observations when contemporaneous or forward weather forecasts are unavailable. Our lag scan ($k \in \{6, 12, 18, 36\}$ steps, 1\,h to 6\,h) demonstrates strictly monotonic cost inflation across all three wind farms. On Hill of Towie, dispatch risk climbs from $0.09022$ at $k=6$ to $0.09257$ at $k=36$, exceeding the unconstrained baseline ($0.08924$) at $k \approx 12$. We emphasize that this experiment characterizes the sensitivity of the training objective to temporal misalignment between the forecast horizon and the historical meteorological proxy, rather than an inference-time telemetry delay (the forecasting backbone takes only past power measurements during inference). Nonetheless, it demonstrates that offline regularization against lagged weather proxies rapidly erodes inductive guidance, presenting an operational limitation for proxy-dependent constraint designs.

\paragraph{Metric Decoupling and the Zero-Bid Diagnostic}
Under standard two-settlement market design, total generator revenue is $R_{\text{total}} = P_{\text{DA}} \cdot q + P_{\text{imb}} \cdot (y - q)$. Evaluations that drop the forward commitment term and report only balancing cashflow are measuring imbalance rent, not forecast skill. As an analytical imbalance-settlement and decision-risk diagnostic, the degenerate Zero-Bid baseline ($q \equiv 0$) provides a mathematical \textit{reductio ad absurdum}: it extracts £223k--£243k per turbine-year across all three farms without generating forecasts, simply by converting all physical generation into uncontracted spill rewarded at $P_{\text{imb}}$. This confirms that standalone unhedged balancing cashflows are decoupled from forecasting skill, and that assessing full commercial value requires joint modeling of forward contract revenue ($P_{\text{DA}} \cdot q$) and imbalance settlement. Operational evaluations must report asymmetric Pinball loss at the market penalty ratio or full two-settlement P\&L.

\subsection{External Validity Boundaries}
\label{sec:boundary_contract}
Five operational boundaries define the validity envelope of our empirical findings:

\paragraph{Boundary 1: Sample Volume and Hypothesis-Space Saturation}
Within-site learning curves confirm that physical constraint benefits concentrate under data-scarce regimes in our tested configurations ($N_{\text{train}} \le 1,000$, Figure~\ref{fig:within_site_learning_curve}). As training data expands to $N_{\text{train}} = 5{,}000$, hypothesis-space saturation reduces the gap between genuine-prior and unconstrained models to within 0.33\% on Kelmarsh, while scrambled priors remain severely detrimental throughout. Continuous corruption tests further establish a narrow fault-tolerance window ($\alpha^* \approx 0.07$ on Penmanshiel, Figure~\ref{fig:corruption_gradient}). Beyond this threshold, engineering drift exceeding 10\% aerodynamic capacity loss or $\pm 5\%$ anemometer bias flips the nominal physical prior from an asset into an operational liability.

\paragraph{Boundary 2: Proxy Misalignment and Observation Staleness}
For proxy-dependent physical constraints, regularization effectiveness depends on temporal alignment between meteorological proxies and target horizons ($k \le 6$ steps). Beyond this, monotonic cost inflation is observed on all three farms ($+2.6\%$ on Hill of Towie at $k=36$). Offline training objectives that substitute stale historical proxies for contemporaneous weather data induce substantial performance degradation.

\paragraph{Boundary 3: Market Clearing and Two-Settlement Structure}
Economic findings serve as an imbalance-settlement and decision-risk diagnostic bounded by two-settlement markets with asymmetric shortage penalties ($r \ge 1.0$) and UK BSC 30-minute discrete netting. Standalone imbalance cashflows do not represent full commercial value and do not generalize to symmetric single-settlement systems or flat feed-in tariffs.

\paragraph{Boundary 4: Computational Graph Audit and Implementation Parity}
Software conclusions apply to PCT commit \texttt{65941ec} and the computational graphs audited here. Unreleased proprietary codebases require independent static graph audit before implementation fidelity claims can be admitted.

\paragraph{Boundary 5: Literature Scoping Scope and Architectural Probe Envelopes}
The systematic literature audit is strictly bounded to the 46 studies identified under our PRISMA-ScR protocol; the documented 95.7\% code deficit and 50.0\% confounded ablation prevalence characterize reporting practices within this scoping cohort rather than an exhaustive census of all hybrid machine learning. Similarly, our PatchTST evaluations function as an architectural sensitivity probe under fixed nominal regularization ($\lambda=0.01$); while demonstrating that qualitative downside risk reduction persists across non-linear deep representations, joint cross-backbone $\lambda$ grid searches remain an open direction for future benchmarking.

\subsection{Two Conjectures Falsified During the Audit}
Two initial hypotheses were refuted during the empirical audit. First, forensic tensor derivation showed that PCT's product-kernel KDE preserves bivariate coupling algebraically, refuting the conjecture of marginal independence. Second, market diagnostic probing showed that post-calibration cashflow inversions are governed by settlement price asymmetries rather than physical regularizers. Documenting these falsification outcomes provides a transparent account of our investigative trajectory and prevents post-hoc rationalization.

\subsection{REFORMS-PCWPF Compliance Protocol}
\label{sec:reforms_pcwpf}
Table~\ref{tab:reforms_checklist} proposes five verification gates as a methodological reporting checklist for physics-constrained wind forecasting studies:

\begin{table*}[!t]
\centering
\footnotesize
\renewcommand{\arraystretch}{0.88}
\setlength{\tabcolsep}{3.5pt}
\caption{Methodological Verification Checklist for Physics-Constrained Wind Power Forecasting (Criteria C1--C5).}
\label{tab:verification_checklist}
\begin{threeparttable}
\begin{tabularx}{\textwidth}{p{1.4cm}p{2.6cm}>{\raggedright\arraybackslash}X>{\raggedright\arraybackslash}Xp{3.2cm}}
\toprule
\textbf{Criterion} & \textbf{Governance Dimension} & \textbf{Mandatory Audit Protocol} & \textbf{Common Pitfall / Anti-Pattern} & \textbf{Required Verification Artifact} \\
\midrule
\textbf{C1} & \textbf{Isolated Prior Ablation} & Train identical network backbones with constraint penalty $\lambda=0$ vs. $\lambda>0$ under identical seed schedules. & Confounding module excision (removing entire GNNs or attention blocks) with physical knowledge. & Table of $\lambda=0$ vs. $\lambda>0$ with identical parameter counts and multi-seed variances. \\
\midrule
\textbf{C2} & \textbf{Domain Specificity Control} & Evaluate pseudo-scrambled priors that preserve loss formulation, architecture, and regularization weight $\lambda$ while permuting physical relations. & Mistaking non-specific numerical regularization under weak penalty settings (e.g., $\lambda = 0.01$ under the present loss scaling) for domain reasoning. & Scrambled control performance curve and empirical specificity ratio ($p$-value contrast). \\
\midrule
\textbf{C3} & \textbf{Proxy Staleness Stress-Test} & If constraints depend on physical inputs (e.g., wind speed), scan historical proxy lag $k \in \{6, 12, 18, 36\}$ steps to evaluate proxy staleness when forward NWP forecasts are absent. & Assuming temporally aligned physical proxies when constructing the offline constraint target. & Observation staleness degradation curve demonstrating operational harm horizon. \\
\midrule
\textbf{C4} & \textbf{Metric Decoupling \& Accounting} & Report symmetric error (RMSE), non-negative asymmetric decision risk (Pinball), and test degenerate zero-bid anomaly baselines. & Reporting unhedged standalone imbalance cashflow in isolation, mistaking accidental price-spike alignment for skill. & Decoupled market evaluation coupling forward commitments and balancing cashflow. \\
\midrule
\textbf{C5} & \textbf{Computational Graph Parity} & Verify that software loss implementations match published mathematical equations in argument direction and normalization. & Argument order flip in asymmetric divergences (e.g., PyTorch \texttt{F.kl\_div}) causing reverse optimization. & Static code audit report and unit test verifying loss tensor gradient directions. \\
\bottomrule
\end{tabularx}
\begin{tablenotes}[flushleft]
\footnotesize
\item \textbf{Notes}: Derived from our empirical audit of 46 scoping studies and adapted from general machine-learning reporting principles \cite{kapoor2024reforms,walsh2021dome}. We propose C1--C5 as an operational reporting checklist for future physics-constrained renewable energy forecasting studies.
\end{tablenotes}
\end{threeparttable}
\end{table*}


\begin{itemize}
\item \textbf{Gate M1 (Parameter-Isolated Ablations):} Parameter-matched comparisons ($\lambda=0$ vs.\ $\lambda>0$) should be required to isolate physical prior contributions, avoiding submodule excisions (e.g., stripping GNN layers) that conflate representational capacity with domain knowledge.
\item \textbf{Gate M2 (Pseudo-Scrambled Negative Controls):} A minimum credible evaluation of physical reasoning claims should include scale-matched permutation controls to separate aerodynamic guidance from non-specific numerical smoothing.
\item \textbf{Gate M3 (Operational Proxy Staleness Stress-Tests):} Multi-step constraint evaluations should avoid assuming zero-latency future weather inputs, reporting performance across $k \in \{6, \dots, 36\}$ steps when forward NWP feeds are absent.
\item \textbf{Gate M4 (Coupled Two-Settlement Evaluation):} Evaluations should avoid standalone unhedged balancing cashflows, reporting instead asymmetric Pinball loss at the market penalty ratio or coupled two-settlement P\&L ($P_{\text{DA}} \cdot q + P_{\text{imb}} \cdot (y - q)$).
\item \textbf{Gate M5 (Static Computational Graph Audit):} Reviewers should expect executable training scripts or verified computational graphs, verifying that software loss execution matches published definitions in argument ordering and reduction semantics.
\end{itemize}
